\chapter{Evidence Tracking and Probe Selection}
\label{ch:method}

An extra coding check can change a prediction, but its result may be unreliable. This chapter explains how the project limits
the influence of one observation and chooses which cases receive a check when the number of checks is restricted. The
simulations make the underlying state and evidence faults known, allowing these decisions to be examined directly.

The final section applies a simpler uncertainty ranking to recorded student activity. It tests whether uncertain predictions
contain a larger share of errors. Together, these methods examine how evidence changes an estimate and where further checking
might be useful.

\section{Bounded Bayesian Updating}

The tracker maintains a probability estimate and revises it after each observation. An ordinary Bayesian update gives more
weight to results that strongly distinguish the two possible states. The modification below limits that weight, because a
faulty check can otherwise cause a large change. The bound controls the size of the update. Whether it improves predictions
is tested in the experiments.

For observation channel $c$, category $z$, and binary state $K_t$, let
$O_c(z,k)=P(Z=z\mid K_t=k,c)$. The state is a binary variable defined by the simulator.
Let $p_t$ be the model's estimated probability of $K_t=1$ before the observation, conditional on the available history.
It need not equal the generating process's conditional probability when the model is wrong. With positive likelihoods,
Bayes' rule gives
\begin{equation}
 p_t^+=\frac{p_tO_c(z,1)}{p_tO_c(z,1)+(1-p_t)O_c(z,0)},
 \qquad
 \frac{p_t^+}{1-p_t^+}=\frac{p_t}{1-p_t}\frac{O_c(z,1)}{O_c(z,0)}.
\end{equation}
Taking logarithms yields the ordinary additive log-odds update. This identity assumes the stated observation likelihoods are
appropriate conditional on the history. Applying several fixed channel likelihoods in succession also assumes the relevant
conditional independence. Correlated evidence can violate that assumption.

The project then modifies the update using fixed trust $0\leq\eta_c\leq1$ and clipping limit $\kappa>0$:

\begin{equation}
  \operatorname{logit}(p_t^+)=\operatorname{logit}(p_t)+
  \eta_c\,\operatorname{clip}\!\left(
    \log\frac{O_c(z,1)}{O_c(z,0)},-\kappa,\kappa
  \right).
\end{equation}

For positive channel likelihoods, an interior prior and non-negative trust, one observation contributes at most $\eta_c\kappa$
to log-odds in absolute value. Across $T$ observations, the sum of the absolute evidence contributions is bounded by
$\sum_{t=1}^{T}\eta_{c_t}\kappa$. Learning and forgetting transitions are separate from this bound.
This is a per-observation bound, not a bound on total movement over a complete episode. Repeated observations can add their
contributions, and the learning and forgetting transitions can also move the estimate. Repeated misleading evidence can
therefore still accumulate substantial influence. Neither calibration nor lower prediction error follows from clipping alone.
The robustness finding is empirical, under the declared simulator.

Equal likelihoods contribute zero. With $\eta_c=1$ and no active clipping, the ordinary Bayesian update is recovered. Otherwise,
the bounded value is a modified belief update and need not be the exact posterior under the original observation model.
The channel-aware variants use category and channel. Their shared input's confidence field is retained but does not multiply
the likelihood again. This is a fixed estimator choice, not learned trust.

The frozen tracker uses $\kappa=2$. Trust weights are 0.50 for public answers, 0.90 for executable probes and 0.40 for
self-explanations. These authored weights limit a single observation's log-odds contribution to 1.0, 1.8 and 0.8 respectively.
They express the study's assumptions about evidence handling and were not estimated from learner data.

The numerical implementation restricts probabilities to $[f_0,1-f_0]$, where $f_0=10^{-6}$, before taking log-odds and
after a non-zero update. The equation above describes the update before this final restriction. For a prior already within
these limits, restricting the result cannot enlarge its movement in log-odds. An input outside those limits is first adjusted.
That adjustment is separate from the bounded evidence contribution. Missing evidence retains the supplied prior.

After the evidence update, the tracker forms the next-turn prior using learning probability $l$ and forgetting probability $f$:
\begin{equation}
 p_{t+1}=p_t^+(1-f)+(1-p_t^+)l.
\end{equation}
This ordering separates evidence about the current state from the subsequent transition. For probabilities in $[0,1]$, the
transition remains in $[0,1]$ as a weighted average. This mechanical property does not validate the transition as a model of
human learning. The implementation applies the numerical probability limits before and after this transition as well.
These transitions apply to the study of bounded updates. The probe selection experiment in
Section~\ref{sec:probe-selection} keeps each case's outcome fixed during selection.

\subsection{Tracker Baselines and Confidence}

The five trackers in the bounded update study differ in how they use an observation. All retain the prior when evidence is missing, before any separate
state transition. The last-observation baseline assigns $1-f_0$ after supporting evidence and $f_0$ after opposing evidence,
where the numerical probability floor is $f_0=10^{-6}$. This deliberately discards the previous estimate.

The hard-BKT baseline uses fixed slip $s=0.15$ and guess $g=0.20$. Its likelihood pair for supporting evidence is
$(1-s,g)$, and for opposing evidence it is $(s,1-g)$, ordered by simulated state $K=1$ and $K=0$.
These parameters are authored rather than fitted to learner records.

For the legacy fractional baseline, let $c\in[0.5,1]$ denote the observation's confidence field, as required by the input contract.
Define $\rho=c$ for supporting
evidence and $\rho=1-c$ for opposing evidence. With fixed $\epsilon=0.15$, the implemented likelihood-like quantities are
\begin{equation}
 \widetilde L_1=(1-\epsilon)^\rho\epsilon^{1-\rho},\qquad
 \widetilde L_0=\epsilon^\rho(1-\epsilon)^{1-\rho}.
\end{equation}
Their contribution to log-odds is therefore
\begin{equation}
 \log\frac{\widetilde L_1}{\widetilde L_0}
 =(2\rho-1)\log\frac{1-\epsilon}{\epsilon}.
\end{equation}
Confidence $c=0.5$ contributes zero. Increasing confidence strengthens the contribution in the category's direction.
This is the frozen transformation, not a calibrated probability of correctness or a claim that confidence specifies a
Jeffrey update. The distinction between evidence strength and a prescribed revised belief is discussed in
\cref{ch:related-work} \parencite{chan2003uncertain}.

The remaining baselines use the channel likelihoods defined above: one applies the ordinary likelihood ratio, while the
bounded variant applies trust and clipping. Neither uses the confidence field. Comparing these trackers thus changes the
observation model as well as, in some comparisons, the update size. Only the channel-aware versus bounded-channel-aware
comparison holds the channel model fixed. Even there, trust weighting and clipping are combined unless separately ablated.

\section{Selecting Probes Within a Fixed Budget}
\label{sec:probe-selection}

Each simulated episode in the probe selection study contains 40 cases. Every ordinary policy sees the same prior estimates and permitted features, then
selects exactly 20 probes before any selected outcome is revealed. The allocation is compulsory, and every selected probe has
the same unit cost. Random selection,
uncertainty-only selection, a plug-in clipped-belief comparator, and the reliability-aware clipped-belief heuristic all spend
the full allocation. Zero or negative scores do not stop selection. Never-probe and always-probe define endpoints. This is
one-step batch selection, without conversational adaptation.

For each case, the simulator draws a starting probability from $\operatorname{Beta}(2,2)$ and clips it to $[0.05,0.95]$.
It then draws $Y_i\sim\operatorname{Bernoulli}(p_i)$ and supplies that same clipped $p_i$ to the policy. Each episode contains
ten families, with one case from each of the four probe classes in every family. Starting
probabilities are calibrated by construction. Held-out environments change evidence reliability, error balance, relevance,
inversion, missingness or dependence. They do not test inaccurate starting probabilities. No learning transition occurs within
a case in this experiment, and the policy cannot inspect the hidden outcome or generating channel rates.

\subsection{Calibration and the Implemented Score}

Sensitivity is the chance that a probe passes when $Y=1$. Specificity is the chance that it fails when $Y=0$. Each has a separate
$\operatorname{Beta}(1,1)$ prior. For $x$ successes in $n$ calibration trials, its posterior is
$\operatorname{Beta}(1+x,1+n-x)$. These trials are simulated from declared channel rates, not collected from students.
Three probe classes use 160 trials per parameter, while the sparse class uses 12. Calibration remains fixed during evaluation.
The dense and sparse high-reliability classes share sensitivity 0.90 and specificity 0.92. The other dense classes use
$(0.78,0.82)$ and $(0.92,0.64)$. Each class contributes ten cases per episode. The shared rates of the first two classes let
the experiment vary calibration sample size while holding their generating reliability constant.

Under symmetric zero-one loss, predicting success incurs expected loss $1-p$, while predicting failure incurs expected loss $p$.
Choosing the smaller gives $R(p)=\min(p,1-p)$ for a coherent probability $p$. This follows the standard posterior expected-loss
decision rule \parencite[Section~5.7, pp.~176--177]{murphy2012machine}.
Let $T_\kappa(p,z;\theta)$ denote the clipped update using sensitivity and specificity $\theta$. The frozen ranking score is
\begin{equation}
 s_i=Q_{0.10,\,\theta\mid D_{\mathrm{cal}}}
 \left[R(p_i)-\mathbb{E}_{Z\mid p_i,\theta}
 R\!\left(T_\kappa(p_i,Z;\theta)\right)\right].
\end{equation}
The implementation uses 8,192 reliability draws with a fixed seed and $\kappa=2$. It takes the lower quantile of these scores
and selects the highest-ranked cases. The plug-in comparator substitutes posterior mean reliability into the same score.
The sample quantile uses linear interpolation. This score averages over pass and fail only. It does not model the probability
of receiving no result. Missing results enter the evaluation environments, where they consume a selected probe and leave its
starting probability unchanged. Thus, the ranking calculation has no explicit penalty for a probe that may return no evidence.
The lower quantile is cautious within the assumed calibration model. It provides no guarantee against failure mechanisms
outside that model.

\subsection{When the Selection Score Overstates the Benefit}

The frozen code and some historical artefacts call these scores EVSI. In this report they are called lower-quantile clipped-risk
scores, because EVSI would suggest a stronger decision-theoretic identity than the implementation has. Clipping generally breaks the correspondence between
the updated belief and the joint model used to average over observations. The score can change when the classification does not.
For $p=0.95$, sensitivity and specificity $0.9$, and $\kappa=2$, passing produces belief $0.992927$ and failing produces
$0.719995$. Table~\ref{tab:acquisition-counterexample} shows the resulting decisions.

\begin{table}[htbp]
  \centering
  \begin{tabular}{@{}lcc@{}}
    \toprule
    Information available & Belief in success & Saved prediction \\
    \midrule
    Before the check & $0.950000$ & Success \\
    Check passes & $0.992927$ & Success \\
    Check fails & $0.719995$ & Success \\
    \bottomrule
  \end{tabular}
  \caption{A counterexample using the frozen clipping bound. The belief changes after either possible check result, but the
  predicted class remains success. The check therefore cannot reduce classification error for this case under the stated model.}
  \label{tab:acquisition-counterexample}
\end{table}

The probability of a pass is $0.86$ and the probability of a failure is $0.14$. The implemented calculation gives

\begin{equation}
  0.05-\left[0.86R(0.992927)+0.14R(0.719995)\right]=0.004717.
\end{equation}

Both branches still predict success, so the actual expected classification error remains $P(Y=0)=0.05$ and its reduction is
zero. A pass moves the belief farther from the threshold, while a failure moves it closer. Weighted by their probabilities,
the risk expression evaluated at those modified beliefs falls below its starting value. The predicted class stays the same
in both branches. This counterexample was found during review after the experiment, using the original clipping bound.
It does not measure how often the issue affected the study's rankings, and is not attributed to the decision-theory reference.

There is a second distinction. Ranking calculates hypothetical updates using each sampled reliability setting. Execution always
updates using the calibration posterior means. The ranking score therefore does not integrate the expected loss of the exact
procedure subsequently executed.

To make the distinction precise, let $\bar\theta$ be the calibration means and
$a_{\bar\theta}(z)=\mathbf{1}\{T_\kappa(p,z;\bar\theta)\geq 0.5\}$. Under generating channel $\theta$, this fixed decision
procedure has expected loss
\begin{equation}
 L_\theta(a_{\bar\theta})=
 \sum_{y\in\{0,1\}}\sum_z P_\theta(Y=y,Z=z\mid p)
 \mathbf{1}\{a_{\bar\theta}(z)\ne y\}.
\end{equation}
Missing observations, where permitted, leave the starting belief unchanged and consume the probe allocation. They must be
included in this sum. The actual expected reduction is $R(p)-L_\theta(a_{\bar\theta})$, assuming the stated joint model.
This is an application of expected-loss evaluation, not a new decision-theory result
\parencite[Section~5.7]{murphy2012machine}. It explains the limitation. It was not substituted into the frozen experiment after
its results were known. The fixed clipped rule can increase expected loss, whereas a Bayes-optimal decision with the option to
ignore an observation cannot do worse in expectation under its own coherent model. That distinction does not provide protection
when the assumed model is wrong.

A decision-aligned calculation would therefore follow four steps for each case:

\begin{enumerate}
  \item Record the classification and expected error before requesting a check.
  \item Enumerate every possible check result, including a missing result where the environment permits one.
  \item Apply the same update and threshold that evaluation will execute. Calculate whether that decision is wrong under
  each possible true outcome.
  \item Average those errors under the stated joint model and subtract the result from the error before the check.
\end{enumerate}

The frozen policy does not implement these four steps as one coherent calculation. They specify the correction required before
a future selector can be described as choosing probes by expected classification-error reduction.

The batch objective introduces a further limitation. With equal probe costs and exactly $k$ selections, taking the largest
$s_i$ maximises the sum of the selected scores. It does not generally maximise a lower quantile of their combined benefit.
For uncertain quantities $X_i$, quantiles do not generally add:
\begin{equation}
 Q_{0.10}\!\left(\sum_{i\in S}X_i\right)
 \ne \sum_{i\in S}Q_{0.10}(X_i)
 \quad\text{in general}.
\end{equation}
Here $S$ is a selected set. Shared uncertainty about channel reliability affects its joint outcome. Even independent quantities
need not satisfy equality. For example, let two independent quantities each equal zero with probability $0.06$ and one otherwise.
Each has lower tenth quantile one, but their sum has lower tenth quantile one, rather than two: the probability of a sum at most
one is $1-0.94^2=0.1164$. This mathematical example is separate from the simulator results.

Consequently, cautious individual rankings provide no established lower-quantile guarantee for the batch's total benefit.
Correcting the individual expected-loss calculation would still leave that distinction to address. The frozen study evaluates
the ranking rule as implemented. Optimal selection under joint reliability uncertainty remains unestablished.

The reference labelled \texttt{oracle} uses true channel rates but retains the assumed update and clipped-belief risk score.
It is a true-channel reference, not an optimal accuracy bound. Historical policy identifiers remain unchanged for reproduction.
The final evaluation measures classification errors against generated outcomes separately from the acquisition score. Those
measurements describe the implemented heuristic, not a correctly integrated value-of-information optimiser.

\subsection{How the Equations Are Implemented}

Table~\ref{tab:acquisition-implementation} maps each calculation to its code. Paths start at
\path{development/src/socratic_tutor/acquisition_study/}.
\begin{table}[htbp]
 \centering
 \begin{tabularx}{\textwidth}{@{}>{\raggedright\arraybackslash}p{0.25\textwidth}>{\raggedright\arraybackslash}X@{}}
 \toprule
 Quantity & Implementation and interpretation \\
 \midrule
 Ranking score & \path{policies.py}: \texttt{reliability\_aware\_evsi\_scores} is the historical function name. It computes a lower quantile of changes in clipped-belief risk. \\
 Executed update & \path{evaluation.py}: selected observations update using calibration means, without reranking. \\
 Measured error & \path{evaluation.py}: final thresholded prediction compared with \texttt{latent\_success}. \\
 Primary interval & \path{primary_analysis.py}: paired episode resampling with a fixed calibration fit. \\
 \bottomrule
 \end{tabularx}
 \caption{Where the probe selection calculations are implemented. The ranking score uses a risk expression evaluated at modified
 beliefs. Measured error counts incorrect final predictions. A higher ranking score need not produce fewer errors.}
 \label{tab:acquisition-implementation}
\end{table}

\subsection{Failure Settings and Their Limits}

\Cref{tab:acquisition-faults} describes the seven stress tests in \path{v1-environments.yaml}.
For sensitivity or specificity $r$, scaling by $s$ gives $0.5+s(r-0.5)$ before inversion.

\begin{table}[htbp]
 \centering
 \begin{tabularx}{\textwidth}{@{}>{\raggedright\arraybackslash}p{0.23\textwidth}>{\raggedright\arraybackslash}X@{}}
 \toprule
 Setting & Constructed change and purpose \\
 \midrule
 Matched & Uses the calibration-generating channel rates to check behaviour when the assumed channel family is appropriate. Finite calibration samples still introduce estimation error. \\
 Lower reliability & Scales both rates' distance from chance by 0.5 to test overconfidence in channel quality. \\
 Asymmetric errors & Scales sensitivity's distance from chance by 0.2 and leaves specificity unchanged to test excess false negatives. \\
 Irrelevant evidence & Sets sensitivity and specificity to 0.5 to test evidence carrying no information about the target. \\
 Inverted evidence & Reverses every sampled result to test a fully reversed observation channel. \\
 Missing results & Missing probability is $0.05+0.55(1-2|p-0.5|)$, capped at 0.60. Lost observations concentrate near uncertain priors. \\
 Family failures & Reverses all four outcomes in a family together with probability 0.30 to test shared corruption. \\
 \bottomrule
 \end{tabularx}
 \caption{Authored acquisition stress settings. The probabilities were chosen to specify controlled experiments. Their values
 do not estimate the frequency of software faults or learner behaviour.}
 \label{tab:acquisition-faults}
\end{table}

The v1 evaluation motivates checking evidence quality, but its observed defects do not establish these probabilities or validate
each mechanism. In particular, full inversion is a deliberately severe control. The missingness variable is proximity to prior
probability one half, not an observed measure of student difficulty.

The broader concern about acquisition under wrong model assumptions is supported by existing work. For example,
\textcite{tang2026robust} analyse error amplification under model misspecification and propose an acquisition criterion that
accounts for it. That work provides related theoretical context. This study does not implement its criterion or transfer its
guarantees to clipped-belief selection.

The matched setting is a project-authored implementation check. Lower reliability and
asymmetric errors are project-authored parameter shifts. Irrelevance and full inversion are project-authored contrastive controls.
Missingness and family corruption are project-authored structural stress conditions. Their immediate specification source is the
frozen configuration, and their purpose is stated in \cref{tab:acquisition-faults}. Neither the cited literature nor v1 supplies
their numerical rates. This distinction avoids presenting plausible failure stories as observed educational data.

Family corruption changes marginal reliability as well as dependence: a base rate $r$ becomes $0.3+0.4r$ after random family
inversion. Without an independent-corruption condition at the same marginal rates, that setting cannot isolate the causal effect
of correlation alone. Its result concerns the combined failure mechanism. No additional control was added after viewing results.

\subsection{Comparison and Uncertainty}

Development checks used 50 episodes per environment. The final comparison used 2,000. The frozen plan permitted one final
evaluation and exact retries, while prohibiting parameter selection from evaluation outcomes or use of those outcomes to
update reliability. Its software-failure rule required invalidating the affected run, fixing the version and rerunning.
Development ablations and calibration sensitivity remain diagnostic analyses, with their own records and denominators.

The acquisition extension had a planned 40-hour technical limit to protect time for writing and verification. This records
the intended scope. The retained experiment reports do not independently establish the total working hours spent. The later
mathematical review corrected interpretation of the saved results without replacing the evaluated method.

A proposed replay of the acquisition policies over the external-model benchmark was dropped. Those records did not contain
the corresponding reliability inputs available before the outcomes were known. Constructing them retrospectively would make
the comparison depend on choices informed by those outcomes. The external benchmark therefore motivates the acquisition
study. It does not provide an external validation of the selector. Replaying saved benchmark responses to check the original
scoring remains a separate operation.

Before the final simulator run, a reviewer with self-described machine-learning experience assessed the design packet and
environment descriptions. The disposition records 6 September 2026. The reviewer reported working independently, without
seeing development outcomes or policy rankings, and recommended proceeding unchanged. They nevertheless marked uncertainty
about omitted failure types, inference of the hidden environment and the episode as an independent unit. These reservations
remain limitations. The recommendation does not resolve them or validate the simulator as a model of learners.

The packet and response were not committed separately, so Git history alone cannot establish their order. The retained
response and disposition document the reported process. Information isolation applies to the fixed, inspected policies.
Arbitrary code can still read local files. Episodes are the resampling unit because failures
may be shared within their case families, while separate episodes use independent random draws. This is an assumption of
the simulation design, rather than an empirical finding from the review.

The primary comparison equally weights six specified mismatch environments, with 2,000 episodes per environment.
Pairing is preserved within each environment. The three primary comparisons use Bonferroni-adjusted percentile-bootstrap
intervals from 20,000 draws. These intervals are conditional on the frozen calibration fit. They do not include variability
from collecting another calibration sample. The separate 20-seed development sensitivity analysis does not change that scope.

The success rule requires favourable aggregate intervals against all three comparators and no adverse per-environment point
estimate against the plug-in comparator. This is the prespecified study criterion. A point estimate cannot rule out meaningful
harm, so the rule is not a per-environment safety guarantee. Equal weighting describes the chosen stress-test mixture, not
deployment prevalence. The candidate failed this rule. The threshold and comparators remained fixed.

\section{Prediction from Historical Learner Records}

The exploratory analysis uses the CSEDM 2019 Data Challenge, version 1.1, which documents novice Python programming records and ten student-separated
evaluation folds \parencite{csedm2019challenge}. The local inventory identifies the analysed archive by its SHA-256 checksum.
The separate Java CodeWorkout release is excluded. Each target's learner-history features use only that learner's earlier events. Preprocessing is fitted within
the training fold. Overall prevalence, problem frequency, and a fixed learner-history logistic model produce out-of-fold
predictions. The primary model uses earlier activity, submission, correctness, hint-request and distinct-problem counts or rates,
plus training-learner success on the target problem. The latter excludes the target's own label for training rows. Test features
use training-fold statistics only. Missing values use training medians and indicators. Numeric scaling is fitted on training data.
Problem success rates use the available training-fold labels without a global calendar cut-off. The evaluation therefore
tests prediction for learners excluded from model fitting on the existing problem set. It does not recreate deployment in
which only records collected before a particular date are available for training.

The chronology boundary is strict: an event contributes features only when its \texttt{Order} is less than the target's
\texttt{StartOrder}. The supplied \texttt{FirstCorrect} value is the outcome, never an input. The source defines the target
as first-attempt success without requesting a hint \parencite[The Data, Predict.csv]{csedm2019challenge}. This study uses
that supplied label rather than reconstructing it from later events. The target's later correctness,
hint use and attempt count are also excluded. Learners could attempt problems in different orders, so problem identity alone
does not establish which events were available at prediction time.

The source documentation notes that activity correctness was reconstructed after the original study and may differ from the
feedback learners received. Version 1.1 also corrected leakage in the supplied example model. This dissertation uses its own
fold-specific pipeline. Reproducing the supplied example's metrics checks agreement with that reference. The documented data
limitations still apply. The local inventory retains unavailable activity correctness values. They are excluded from
correctness-rate denominators while activity counts retain the events.

The predictor uses scikit-learn \parencite{pedregosa2011scikit}. The locked environment specifies the installed version.
The frozen \path{v1-analysis.yaml} CSEDM plan specifies logistic regression with an intercept, L2 regularisation and inverse
regularisation strength $C=1$. Training uses the \texttt{lbfgs} solver, no class weighting, at most 1,000 iterations, and tolerance
$10^{-8}$. The recorded random seed is 9052901. The implementation treats a convergence warning as a failed run rather than
silently accepting an unfinished fit. These settings are fixed. There is no hyperparameter search, fitted probability calibration,
or model selection based on test-fold results.

Probabilities at or above 0.5 predict first-attempt success. The separate uncertainty ranking uses $|p-0.5|$ in ascending order.
Small values mean the classifier is near its decision threshold. This is a simple ranking score, not a full estimate of model
uncertainty or a guarantee that a probability is calibrated. Equal scores are ordered by a stable hash of the fold, learner,
problem and prediction cut-off identifiers, without using the outcome.

For fold $f$ with $N_f$ targets, selection takes $k_f=\lfloor0.5N_f\rfloor$ predictions closest to probability $0.5$.
This selects 361 of 729 targets, rather than the 364 from a single global allocation. If $E_f$ is the fold's error count, expected
random error capture is $\sum_f k_f E_f/N_f$. Dividing this by $\sum_f E_f$ gives the random captured-error recall used in the
primary comparison. The observed recall uses the actual number of selected errors instead.

Accuracy, Brier score, log loss and calibration error are secondary descriptions. Log loss clips probabilities at $10^{-6}$
and $1-10^{-6}$. Calibration error uses ten equal-width bins. Leave-one-fold-out summaries check whether omitting a single fold
changes the primary pattern. They do not retrain models or create independent replications.

The 10,000-draw bootstrap resamples learners, retaining all their rows and the original predictions and selections. Its interval
is conditional on those fitted predictions and selections. Models are not refitted and cases are not reranked. The original
fold selection probabilities $k_f/N_f$ are also held fixed when calculating the resampled random comparator. Shared training
sets across folds create further dependence that learner resampling does not fully represent. The analysis concerns the existing
learners and problems. It establishes neither unseen-problem generalisation nor whether reviewing an error would correct it.
